\documentclass[12pt]{article}
\usepackage[utf8]{inputenc}
\usepackage[T1]{fontenc}
\usepackage[margin=1in]{geometry}
\usepackage{setspace}
\usepackage{xcolor}
\usepackage[colorlinks=true, urlcolor=blue!60!black, linkcolor=blue!60!black]{hyperref}
\setlength{\parindent}{2em}
\setlength{\parskip}{1.25em}
\renewcommand{\baselinestretch}{1.0}
\begin{document}
\begin{center}
{\LARGE \textit{City of Cleburne v. Cleburne Living Center, Inc.}, 473 U.S. 432 (1985)\par}
\vspace{0.8em}
{\large Case Brief --- Four-Page Edition\par}
\end{center}
\vspace{1em}

\section{Memory Jogger}

Intellectual disability is not a quasi-suspect classification triggering heightened Equal Protection scrutiny, yet a city violates rational basis review when it denies a special use permit for a group home based on irrational prejudice and unsubstantiated neighborhood fears (\href{https://supreme.justia.com/cases/federal/us/473/432/}{473 U.S. 432}).

\section{Detailed Case Facts}

In 1980 Jan Hannah bought a four-bedroom house at 201 Featherston Street in Cleburne, Texas, to lease to Cleburne Living Center (CLC) as a group home for thirteen men and women with intellectual disabilities, operating as a closely regulated, federally reimbursed intermediate care facility. The property lay in an ``R-3'' district that permitted by right apartment houses, boarding houses, fraternity houses, dormitories, hospitals, and nursing homes, but required an annually renewed special use permit for ``[h]ospitals for the insane or feeble-minded, or alcoholic [sic] or drug addicts, or penal or correctional institutions.''

The city classified the home as a ``hospital for the feebleminded,'' and the City Council voted 3 to 1 to deny a permit. Its reasons were the negative attitudes and fears of nearby owners and elderly neighbors, possible harassment by students at a junior high school across the street (which itself enrolled about 30 students with intellectual disabilities), the 500-year flood plain, legal responsibility for residents' actions, and the home's size and occupancy. The dispute arose amid the deinstitutionalization movement and new federal and state protections for persons with disabilities.

\section{Procedural History}

CLC, Hannah, and Advocacy, Inc.\ sued in the Northern District of Texas. The district court found that the home would have been permitted had its residents not been intellectually disabled, but upheld the ordinance under minimum rational basis review. The Fifth Circuit reversed, holding intellectual disability quasi-suspect and invalidating the ordinance on its face and as applied. The Supreme Court affirmed the as-applied invalidation, vacated the rest, and remanded on July 1, 1985.

\section{Judicial Votes}

Affirmed in part and vacated in part; \textbf{9--0} on invalidating the permit denial, \textbf{6--3} on the standard of review.
\begin{itemize}
\item \textbf{Majority:} Justice White, joined by Chief Justice Burger and Justices Powell, Rehnquist, Stevens, and O'Connor.
\item \textbf{Concurrence:} Justice Stevens, joined by Chief Justice Burger.
\item \textbf{Concurring in the judgment in part and dissenting in part:} Justice Marshall, joined by Justices Brennan and Blackmun.
\end{itemize}

\section{Holding}

Classifications based on intellectual disability are not quasi-suspect and receive rational basis review. Even so, a classification must be rationally related to a legitimate state interest, and ``mere negative attitudes, or fear, unsubstantiated by factors which are properly cognizable in a zoning proceeding,'' and ``irrational prejudice'' are not legitimate interests. The permit requirement is invalid as applied to the Featherston home; the Court did not reach facial validity.

\section{Analysis of Opinions}

\textbf{Majority (White).} White gave four reasons, drawing on \textit{Murgia}, for denying quasi-suspect status: persons with intellectual disabilities have ``a reduced ability to cope with and function in the everyday world,'' making their differences relevant to legitimate goals; extensive remedial legislation shows lawmakers addressing their needs rather than acting from antipathy; that responsiveness negates political powerlessness; and otherwise courts could not deny heightened scrutiny to ``the aging, the disabled, the mentally ill, and the infirm.'' But rational basis review still forbids arbitrary classifications and those resting on ``a bare . . . desire to harm a politically unpopular group.'' Because the R-3 zone freely admitted boarding houses, dormitories, hospitals, and nursing homes, the city had to show a special threat from the Featherston home. Neighbor fear cannot be given legal effect (\textit{Palmore v. Sidoti}); the harassment worry was contradicted by the school's own students; and the flood, liability, density, and fire concerns applied equally to permitted uses.

\textbf{Stevens (concurring).} Rejecting rigid tiers, Stevens saw ``a continuum of judgmental responses'' governed by whether ``an impartial lawmaker could logically believe that the classification would serve a legitimate public purpose that transcends the harm to the members of the disadvantaged class.'' Cleburne's permit rested on irrational fears.

\textbf{Marshall (concurring in part, dissenting in part).} Marshall criticized the ``two for the price of one'' rejection of heightened scrutiny and exposed the majority's unacknowledged ``second order'' rational basis review, which departs from \textit{Williamson v. Lee Optical}. Citing a ``lengthy and tragic history'' of segregation, eugenic sterilization, and exclusion, he would apply heightened scrutiny and invalidate the ordinance on its face.

\section{Examples: Future Applications}

\begin{enumerate}
\item \textbf{Same-Side (Invalid): Autism Group Home.} A suburb permits group residences for students or cardiac patients by right but requires a permit and neighbor consent for a supervised home for eight adults with autism. Its traffic and property-value concerns apply identically to permitted dwellings.
\item \textbf{Same-Side (Invalid): Cosmetology Ban.} A state bars anyone with a history of depression or bipolar disorder from cosmetology licenses while licensing others with chronic conditions after testing. The exclusion rests on stereotypes, not capacity.
\item \textbf{Opposite-Side (Upheld): Bus-Driver Testing.} A cognitive-reaction exam for commercial bus drivers disqualifies applicants with moderate intellectual disabilities. Highway safety is rationally served by objective testing.
\item \textbf{Opposite-Side (Upheld): Neutral Occupancy Cap.} An R-1 ordinance limits households to four unrelated people, incidentally blocking an eight-resident group home. As Footnote 8 noted, citing \textit{Belle Terre}, a neutral cap does not single out disability.
\item \textbf{Fence-Sitter: Residential Care Permits.} A city requires permits for all residential care facilities over eight residents, citing evacuation times. It treats comparable populations alike, but if equally large dormitories are exempt without proof of higher risk, the line may rest on stereotypes.
\end{enumerate}

\section{Critique}

\textbf{Doctrinal and progressive views.} Scholars dubbed the approach ``rational basis with bite'': White demanded record evidence and condemned underinclusiveness under a label that promises deference, leaving lower courts without criteria for when searching review begins, a tension resurfacing in \textit{Romer} and \textit{Windsor}. Echoing Marshall, antisubordination critics note the paradox that legislative progress stripped the group of heightened protection, with consequences in \textit{Garrett} (2001), and fault the remedy for leaving the ``feeble-minded'' text intact.

\textbf{Originalist and textualist views.} Restraint-minded formalists argue the Court second-guessed local zoning and inferred prejudice without textual mooring. Libertarian originalists such as Joseph Diedrich counter that the Fourteenth Amendment requires courts to evaluate actual ends and evidence, so the inquiry was sound but should have been candid.

\section{Key Quotations}

\begin{quote}
``The Equal Protection Clause of the Fourteenth Amendment commands that no State shall `deny to any person within its jurisdiction the equal protection of the laws,' which is essentially a direction that all persons similarly situated should be treated alike.''
\end{quote}
\noindent --- Justice White, for the Court, 473 U.S. at 439.

\begin{quote}
``But mere negative attitudes, or fear, unsubstantiated by factors which are properly cognizable in a zoning proceeding, are not permissible bases for treating a home for the mentally retarded differently from apartment houses, multiple dwellings, and the like.''
\end{quote}
\noindent --- Justice White, for the Court, 473 U.S. at 448.

\begin{quote}
``The short of it is that requiring the permit in this case appears to us to rest on an irrational prejudice against the mentally retarded, including those who would occupy the Featherston facility and who would live under the closely supervised and highly regulated conditions expressly provided for by state and federal law.''
\end{quote}
\noindent --- Justice White, for the Court, 473 U.S. at 450.

\end{document}
